% ============================================================
%  style/schlagmuster.tex — Schlagmuster rechts am Rand
%  (wird von style.tex eingebunden)
% ============================================================
% ============================================================
%  Schlagmuster (rechts am Rand, über den Text gelegt; braucht zwei
%  Durchläufe, latexmk erledigt das) & Griffe (normal im Textfluss)
% ============================================================
\usepackage{tikzpagenodes}
\usepackage{svg} % nur für reine .svg-Dateien (braucht Inkscape + --shell-escape)

% Einheitlicher Maßstab für alle Schlagmuster (gleiche Strichstärke/Notengröße);
% 0.32 -> ein 16tel-Takt aus Inkscape (~352pt breit) wird ca. 40mm breit
\newcommand{\schlagmusterskala}{0.32}
\newlength{\griffeabstand}\setlength{\griffeabstand}{3mm}  % Abstand zw. Griffen
\newlength{\griffeoben}\setlength{\griffeoben}{3mm}       % Abstand zum Text darüber
\newcommand{\schlagmusterdeckkraft}{0.44} % 0.44 * Schwarz auf Papier ~ inkmeta (#8e8e8e)
\newcommand{\schlagmusterfont}{\serif\itshape\fontsize{7.5}{8}\selectfont\color{inkmeta}}
\newsavebox{\overlaybox}

% Bild aus graphics/: bevorzugt PDF, sonst SVG
\newcommand{\schlagmusterbild}[2]{%
  \IfFileExists{graphics/#2.pdf}{\includegraphics[scale=#1]{graphics/#2.pdf}}{%
    \IfFileExists{graphics/#2.svg}{\includesvg[scale=#1]{graphics/#2}}{%
      \PackageWarning{style}{Schlagmuster 'graphics/#2' nicht gefunden}%
      \fbox{\ttfamily #2?}}}%
}

% \schlagmuster[Skalierung]{Datei ohne Endung}[Beschriftung]
% Skalierung optional (Standard \schlagmusterskala) -> alle Muster gleich groß,
% die Breite ergibt sich aus der Länge des Musters.
% Erscheint am rechten Textrand auf Höhe der Zeile, in der es steht.
% In Songs einfach hinten an eine Zeile hängen, z.B.
%   \slbl{Chorus:}\schlagmuster{summer-of-69-v}
\NewDocumentCommand{\schlagmuster}{ O{\schlagmusterskala} m o }{%
  \edef\schlagmusterdatei{\detokenize{#2}}% alltt macht "-" aktiv
  \sbox{\overlaybox}{\schlagmusterbild{#1}{\schlagmusterdatei}}%
  % Fläche als Hindernis für \griffe merken (ganze Breite rechts blockiert)
  \griffschreiben{1073741823}%
    {\number\dimexpr0.8\baselineskip\IfNoValueF{#3}{+4mm}\relax}%
    {\number\dimexpr\ht\overlaybox+\dp\overlaybox-0.8\baselineskip\relax}%
  \begin{tikzpicture}[remember picture,overlay]
    % schwarzes Bild halbtransparent -> wirkt im selben Grau wie inkmeta
    \node[anchor=north east,inner sep=0pt,opacity=\schlagmusterdeckkraft] (sm)
      at ($(current page text area.east |- 0,0)+(0,0.8\baselineskip)$)
      {\usebox{\overlaybox}};
    \IfNoValueF{#3}{\node[anchor=south east,inner sep=0pt] at ($(sm.north east)+(0,0.5mm)$)
      {\schlagmusterfont #3};}
  \end{tikzpicture}%
}

